```latex
\documentclass[11pt]{article}

\usepackage{amsmath,amssymb,amsthm}
\usepackage{hyperref}

\newtheorem{theorem}{Theorem}
\newtheorem{proposition}[theorem]{Proposition}
\newtheorem{lemma}[theorem]{Lemma}
\newtheorem{remark}[theorem]{Remark}

\newcommand{\Mg}{\mathcal M_g}
\newcommand{\Mgn}[2]{\mathcal M_{#1,#2}}
\newcommand{\Mbar}{\overline{\mathcal M}_g}
\newcommand{\C}{\mathbb C}
\newcommand{\D}{\Delta}

\begin{document}

\begin{center}
{\Large\bf On ``non-varying factors'' in boundary strata: a caveat and a correction}\\[0.5em]
\end{center}

\begin{abstract}
We explain that the conclusion requested in the prompt (constancy of the point $Y_t$ on a non-varying factor) does \emph{not} follow from the stated stratified-triviality assumptions alone: one can construct holomorphic families $V_t\subset \mathcal M_g$ whose boundary curves in a fixed stratum $\Delta_\Gamma$ are of the form $\{Y_t\}\times C$ with $Y_t$ varying holomorphically. We then record a standard sufficient strengthening (namely that $V_t$ remain totally geodesic) under which the desired conclusion holds, via finiteness of higher-dimensional totally geodesic subvarieties.
\end{abstract}

\section{Setup and the claim}

Let $\Mg$ be the moduli space of smooth genus $g$ curves, and $\Mbar$ its Deligne--Mumford compactification with boundary stratification by dual graphs:
\[
\Mbar=\Mg\sqcup \partial\Mg,\qquad
\partial\Mg=\bigsqcup_G \Delta_G.
\]
Fix a dual graph $\Gamma$ such that the stratum $\Delta_\Gamma$ is (as in the prompt) a product
\[
\Delta_\Gamma \cong \Mgn{g_1}{n_1}\times\cdots\times \Mgn{g_k}{n_k},
\]
and denote by $\pi_i$ the projection to the $i$-th factor.

Let $V_0\subset \Mg$ be a smooth complex surface and $\{V_t\}_{t\in\D}$ a holomorphic family of smooth complex surfaces deforming $V_0$, with the technical assumptions (i)--(ii) in the prompt about smoothly trivial families on each boundary stratum and compatibility under passing to deeper strata.

For each $t$, define
\[
\partial V_t:=\overline{V_t}\cap \Delta_\Gamma,
\]
where the closure is taken in $\Mbar$.

The prompt asks one to prove:

\medskip
\noindent\textbf{Claim (as stated).}
\emph{Assume $\dim_\C(\partial V_t)=1$ and there exists an index $i$ such that $\pi_i(\partial V_t)=\{Y_t\}$ is a single point for all $t\in\D$. Then $Y_t$ is independent of $t$, i.e. $Y_t=Y_0$ for all $t$.}
\medskip

\section{The claim is false under the stated hypotheses (counterexample)}

The key point is that the condition ``$\partial V_t$ is constant on the $i$-th factor'' only says that each boundary curve lies in a \emph{fiber} of $\pi_i$, but it does not prevent those fibers from moving in a holomorphic family. This is not ruled out by (i)--(ii).

We now give a concrete algebro-geometric construction inside moduli space using standard plumbing coordinates.

\subsection{Plumbing around a separating boundary divisor}

Fix a decomposition $g=g_1+g_2$ with $g_1,g_2\ge 1$, and let $\Gamma$ be the dual graph of a stable curve with one separating node, whose normalization has two components of genera $g_1$ and $g_2$ with one marked point each (the preimage of the node). Then
\[
\Delta_\Gamma \cong \Mgn{g_1}{1}\times \Mgn{g_2}{1}.
\]
(We ignore finite-quotient issues as in the prompt.)

The Deligne--Mumford stack $\overline{\mathcal M}_g$ (and hence its coarse moduli space locally) is locally analytically a normal crossings compactification; in particular, near a point of this boundary divisor, one has holomorphic \emph{plumbing} (or \emph{smoothing}) coordinates in which a neighborhood of $\Delta_\Gamma$ is modeled by
\[
U_1\times U_2 \times \D_q \;\longrightarrow\; \Mbar
\]
where $U_j\subset \Mgn{g_j}{1}$ are small analytic neighborhoods and $q$ is the smoothing parameter of the node, with $q=0$ corresponding to the nodal curve and $q\neq 0$ giving a smooth curve in $\Mg$; see e.g.~\cite[Ch.~3]{HM} or the original construction in \cite{DM}.

Concretely, we obtain a holomorphic map
\[
\mathrm{pl}:U_1\times U_2 \times \D \to \Mbar
\]
such that:
\begin{itemize}
\item $\mathrm{pl}(x_1,x_2,0)\in \Delta_\Gamma$ corresponds to the nodal union of the pointed curves $x_1\in U_1$, $x_2\in U_2$ glued at their marked points;
\item $\mathrm{pl}(x_1,x_2,q)\in \Mg$ for $q\neq 0$ (smoothing the node).
\end{itemize}

\subsection{A family of surfaces with a moving non-varying factor in the boundary}

Fix:
\begin{itemize}
\item a nonconstant holomorphic map $a:\D\to U_1\subset \Mgn{g_1}{1}$, $t\mapsto a(t)$;
\item a (connected) smooth complex curve $C\subset U_2\subset \Mgn{g_2}{1}$.
\end{itemize}
For example, one may take $C$ a small analytic disk in $U_2$; if one insists on algebraicity one may choose $C$ to be an algebraic curve in $\Mgn{g_2}{1}$.

Define, for each $t\in\D$, the surface
\[
V_t \;:=\; \mathrm{pl}\bigl(\{a(t)\}\times C \times (\D^\ast)\bigr)\;\subset\; \Mg,
\]
where $\D^\ast=\D\setminus\{0\}$.

\begin{lemma}\label{lem:Vt-smooth}
For $U_1,U_2,\D$ chosen sufficiently small, each $V_t$ is a smooth complex surface in $\Mg$, and the union
\[
\mathcal V:=\bigcup_{t\in\D} (V_t\times\{t\}) \subset \Mg\times \D
\]
is a complex submanifold defining a holomorphic family deforming $V_0$.
\end{lemma}

\begin{proof}
Consider the holomorphic map
\[
F:\D_t\times C \times \D_q^\ast \to \Mg\times \D_t,\qquad
F(t,b,q) :=\bigl(\mathrm{pl}(a(t),b,q),\,t\bigr).
\]
For $U_1,U_2,\D$ small, $\mathrm{pl}$ is a holomorphic immersion in the $(b,q)$-variables (this is part of the local product/normal-crossings structure of $\Mbar$ along $\Delta_\Gamma$; see \cite[Ch.~3]{HM}), hence $F$ is a holomorphic immersion. Its image is precisely $\mathcal V$, and the fiber over fixed $t$ is $V_t$. Smoothness of each $V_t$ follows.
\end{proof}

Now consider the closure $\overline{V_t}\subset \Mbar$; by construction it is obtained by allowing $q\to 0$ in the plumbing family. Therefore
\[
\partial V_t=\overline{V_t}\cap \Delta_\Gamma = \{a(t)\}\times C \subset \Mgn{g_1}{1}\times \Mgn{g_2}{1}.
\]
In particular $\dim_\C(\partial V_t)=1$, and if we take $i=1$ then
\[
\pi_1(\partial V_t)=\{a(t)\}=: \{Y_t\}
\]
is a single point for every $t$, \emph{but $Y_t$ varies with $t$} since $a$ is nonconstant.

\begin{proposition}\label{prop:counterexample}
The family $\{V_t\}$ constructed above satisfies the hypotheses in the prompt concerning holomorphicity and stratified smooth triviality along $\Delta_\Gamma$, and it yields a family $\partial V_t$ with $\pi_1(\partial V_t)=\{Y_t\}$ varying with $t$. Hence the conclusion $Y_t\equiv Y_0$ does \emph{not} follow from the stated assumptions alone.
\end{proposition}

\begin{proof}
We have already verified holomorphicity and smoothness in Lemma~\ref{lem:Vt-smooth}. The description
\[
\partial V_t=\{a(t)\}\times C
\]
shows explicitly that $\partial V_t$ varies smoothly (indeed holomorphically) with $t$ inside the stratum $\Delta_\Gamma$, and the compatibility with further strata is automatic in this local plumbing model: any deeper boundary stratum inside $\overline{\Delta_\Gamma}$ is obtained by degenerating the pointed curves in $U_1$ or $U_2$, and the intersection loci are products with the same dependence on $t$ via $a(t)$.

Finally, $t\mapsto Y_t=a(t)$ is nonconstant by choice, so $Y_t\not\equiv Y_0$.
\end{proof}

\begin{remark}
The example shows that, from a purely algebro-geometric standpoint of local normal-crossings structure near $\Delta_\Gamma$, there is no obstruction to moving a non-varying factor in a holomorphic family. Therefore, any valid proof of the claim as stated must use additional global/geometric input beyond (i)--(ii), such as strong rigidity properties coming from the Teichm\"uller metric.
\end{remark}

\section{A standard sufficient strengthening under which the desired conclusion holds}

One natural strengthening is to assume the entire family remains totally geodesic.

\begin{theorem}[Constancy under totally geodesic deformation]\label{thm:tg-family}
Assume, in addition to the hypotheses in the prompt, that \emph{each} $V_t$ is totally geodesic (for the Teichm\"uller metric). Then the family is locally constant: after shrinking $\D$ if necessary, we have $V_t=V_0$ for all $t$, and in particular $\pi_i(\partial V_t)=\pi_i(\partial V_0)=\{Y_0\}$ for all $t$.
\end{theorem}

\begin{proof}
By a theorem of Wright, for fixed $(g,n)$ there are only finitely many closed totally geodesic subvarieties of $\mathcal M_{g,n}$ of complex dimension $>1$; equivalently, there are only finitely many totally geodesic subvarieties of $\Mg$ of complex dimension $2$ up to equality in $\Mg$ \cite{WrightTG} (see the abstract and the finiteness statement therein). Therefore a holomorphic map from a connected complex base (here the disk $\D$) into this finite set must be constant. Hence $V_t=V_0$ for all $t$ (after possibly shrinking $\D$ to ensure we remain in the regime of the finiteness theorem), and the asserted constancy of $\pi_i(\partial V_t)$ follows.
\end{proof}

\begin{remark}
Theorem~\ref{thm:tg-family} is \emph{not} what the prompt asks (it assumes total geodesicity for all $t$), but it illustrates the kind of additional rigidity input that is needed to force constancy of $Y_t$.
\end{remark}

\section{Conclusion}

Proposition~\ref{prop:counterexample} shows that the requested conclusion $Y_t\equiv Y_0$ does not follow from the stratified smooth triviality and compatibility assumptions alone. Any correct proof must employ stronger hypotheses tying the deformation to the Teichm\"uller-geodesic geometry (or other rigidity phenomena). One such strengthening is given in Theorem~\ref{thm:tg-family}.

\begin{thebibliography}{99}

\bibitem{DM}
P.~Deligne and D.~Mumford,
\emph{The irreducibility of the space of curves of given genus},
Publ.\ Math.\ IH\'ES \textbf{36} (1969), 75--109.

\bibitem{HM}
J.~Harris and I.~Morrison,
\emph{Moduli of Curves},
Graduate Texts in Mathematics, Vol.~187, Springer, 1998.

\bibitem{WrightTG}
A.~Wright,
\emph{Totally geodesic submanifolds of Teichm\"uller space},
arXiv:1702.03249; published in \emph{J.\ Differential Geom.} (2020).

\end{thebibliography}

\end{document}
```